\begin{table}[t]
  \centering
  \caption{\textbf{段階ごとの権威と失敗時の挙動。}
  どの段階も、欠けた権威を下流で静かに補わない。
  }
  \label{tab:component_contracts}
  \small
  \begin{tabularx}{\linewidth}{@{}p{0.16\linewidth}p{0.30\linewidth}X@{}}
    \toprule
    段階 & 公開する権威 & 推測せず棄却するもの \\
    \midrule
    再構成 & シーンスキーマ、復元カメラ、内部パラメータ、
      正規化点群、レンダ可能な 3DGS &
      隠れた再構成経路、壊れた配列、非直交回転、公開コマンドの欠落 \\
    アライメント & 共通メートル尺度、相互のシーン--コート変換、
      トポロジー、fit/holdout 指標 &
      コート一面への縮退、holdout の再当てはめ、
      尺度不一致、意味セグメント被覆の不足 \\
    データ生成 & RGB と深度、対象コート、カメラ変換、
      \KPName の物理番号、幾何妥当性と描画支持、軌道グループ分割 &
      スキーマ変換、最近傍による点の同一視、
      フレーム単位の分割漏洩、視点の差し替え \\
    検出器 & 型付き密教師、姿勢ベクトル、補正後の内部パラメータ、
      有効領域マスク &
      未対応 source や未対応補正、単一ピーク前提での多峰キーポイント、
      姿勢枝またはキーポイント枝の欠落 \\
    \bottomrule
  \end{tabularx}
\end{table}
